%% notes-tex/database/database-expansion-hosts/sections/summary.tex
%% [[experiments.database.expansion-hosts]]
%% https://github.com/Mjvolk3/torchcell/tree/main/notes-tex/database/database-expansion-hosts/sections/summary.tex
%% SOURCE: hand-authored prose around a generated table; tables/summary.tex names its own
%% generating script. The time-axis figures come from the triage report
%% /scratch/projects/torchcell-scratch/tmp/deep-research/reports/Pk and next dataset wave triage.md
%% and its research notes; the schema fields from torchcell/datamodels/schema.py.

\clearpage
\section{Summary statistics\secstatus{todo}}\label{sec:summary}

%% GENERATED FILE -- do not hand-edit.
%% SOURCE: experiments/database/scripts/build_hosts_candidate_datasets_table.py
\begingroup
\small
\begin{longtable}{@{}l r r r r r@{}}
\caption[]{Summary of the recommended set, split at the fifty line. \emph{Top 50} and
\emph{Extra} count rows; \emph{Genotypes}, \emph{Instances} and \emph{Meas.} sum the row
axes over all 60 and are lower bounds, since a row with no count contributes
nothing. The set names 22 distinct schema classes the substrate does not hold
(Table~\ref{tab:schema}).}
\label{tab:summary}\\
\toprule
 & Top 50 & Extra & Genotypes & Instances & Meas. \\
\endfirsthead
\multicolumn{6}{@{}l}{\footnotesize\emph{Table~\ref{tab:summary}, continued}}\\
\toprule
 & Top 50 & Extra & Genotypes & Instances & Meas. \\
\endhead
\bottomrule
\endfoot
\midrule
\multicolumn{6}{@{}l}{\emph{By band}}\\
regulatory DNA & 34 & 1 & 166,797,405 & $1.7\times 10^{8}$ & $2.1\times 10^{8}$ \\
cross-host fitness & 0 & 9 & 0 & $1.4\times 10^{3}$ & $6.2\times 10^{6}$ \\
host bootstrap & 8 & 0 & 260 & $9.6\times 10^{3}$ & $3.5\times 10^{5}$ \\
scale & 8 & 0 & 5,467 & $5.8\times 10^{3}$ & $4.4\times 10^{6}$ \\
\midrule
\multicolumn{6}{@{}l}{\emph{By tier}}\\
tier 1 & 15 & 5 & 166,022,613 & $1.7\times 10^{8}$ & $1.9\times 10^{8}$ \\
tier 2 & 12 & 2 & 757,235 & $1.2\times 10^{6}$ & $3.1\times 10^{7}$ \\
tier 3 & 13 & 1 & 21,139 & $3.0\times 10^{4}$ & $7.7\times 10^{5}$ \\
tier 4 & 10 & 2 & 2,145 & $2.2\times 10^{3}$ & $1.5\times 10^{5}$ \\
\midrule
\multicolumn{6}{@{}l}{\emph{By host}}\\
E. coli & 9 & 0 & 2,861,834 & $5.2\times 10^{6}$ & $4.5\times 10^{7}$ \\
P. aeruginosa & 0 & 1 & 0 & 8 & $3.6\times 10^{4}$ \\
P. fluorescens & 0 & 4 & 0 & 745 & $3.4\times 10^{6}$ \\
P. kudriavzevii & 8 & 0 & 260 & $9.6\times 10^{3}$ & $3.5\times 10^{5}$ \\
P. putida & 8 & 0 & 5,467 & $5.8\times 10^{3}$ & $4.4\times 10^{6}$ \\
P. simiae & 0 & 1 & 0 & 201 & $9.1\times 10^{5}$ \\
P. syringae & 0 & 2 & 0 & 69 & $3.1\times 10^{5}$ \\
S. cerevisiae & 25 & 1 & 163,935,571 & $1.6\times 10^{8}$ & $1.7\times 10^{8}$ \\
S. stutzeri & 0 & 1 & 0 & 350 & $1.6\times 10^{6}$ \\
\midrule
\multicolumn{6}{@{}l}{\emph{By class}}\\
CRISPR library screen & 1 & 0 & 47 & 188 & 188 \\
Deep mutational scan & 3 & 0 & 4,176 & $4.2\times 10^{3}$ & $4.2\times 10^{3}$ \\
Engineering parts & 3 & 0 & 3,100 & $3.1\times 10^{3}$ & $3.1\times 10^{3}$ \\
Fitness / chemical genomics & 0 & 9 & 0 & $1.4\times 10^{3}$ & $6.2\times 10^{6}$ \\
Natural variation & 4 & 0 & 900 & $1.5\times 10^{4}$ & $1.5\times 10^{4}$ \\
Proteome & 1 & 0 & 0 & 7 & $3.9\times 10^{4}$ \\
Regulatory DNA & 27 & 1 & 166,789,446 & $1.7\times 10^{8}$ & $2.1\times 10^{8}$ \\
Transcriptome & 7 & 0 & 20 & 472 & $2.6\times 10^{6}$ \\
Transposon fitness & 4 & 0 & 5,443 & 430 & $2.1\times 10^{6}$ \\
\midrule
\multicolumn{6}{@{}l}{\emph{By sequence basis}}\\
KT2440+transposon & 4 & 0 & 5,443 & 430 & $2.1\times 10^{6}$ \\
KT2440-reference & 1 & 0 & 0 & 321 & $1.8\times 10^{6}$ \\
MG1655+reporter-locus & 2 & 0 & 329,392 & $8.3\times 10^{5}$ & $8.6\times 10^{5}$ \\
Pk-isolate-WGS & 5 & 0 & 210 & $9.4\times 10^{3}$ & $2.9\times 10^{5}$ \\
S288C+designed-edit & 5 & 0 & 503,524 & $5.1\times 10^{5}$ & $5.1\times 10^{5}$ \\
S288C+integration-locus & 4 & 0 & 2,648 & $2.6\times 10^{3}$ & $1.5\times 10^{6}$ \\
SD108+guide & 1 & 0 & 47 & 188 & 188 \\
engineered-chassis & 3 & 0 & 3 & 93 & $5.1\times 10^{5}$ \\
engineered-chassis+guide & 1 & 0 & 0 & 7 & $3.9\times 10^{4}$ \\
fitbrowser+transposon & 0 & 9 & 0 & $1.4\times 10^{3}$ & $6.2\times 10^{6}$ \\
isolate-WGS & 1 & 0 & 683 & 683 & 683 \\
plasmid-reporter-library & 22 & 1 & 165,961,158 & $1.7\times 10^{8}$ & $2.1\times 10^{8}$ \\
pseudomonad-isolate-WGS & 1 & 0 & 24 & $4.9\times 10^{3}$ & $4.9\times 10^{3}$ \\
\midrule
\multicolumn{6}{@{}l}{\emph{By retrieval class}}\\
retrievable programmatically & 30 & 9 & 166,762,265 & $1.7\times 10^{8}$ & $2.2\times 10^{8}$ \\
needs a human with a browser & 17 & 1 & 40,635 & $4.9\times 10^{4}$ & $1.8\times 10^{6}$ \\
needs an ask to a named party & 3 & 0 & 232 & $1.4\times 10^{4}$ & $1.4\times 10^{4}$ \\
\midrule
\multicolumn{6}{@{}l}{\emph{By verification label}}\\
counts read on the source & 34 & 0 & 166,787,537 & $1.7\times 10^{8}$ & $2.1\times 10^{8}$ \\
counts from a search snippet & 7 & 1 & 4,421 & $1.3\times 10^{4}$ & $1.6\times 10^{5}$ \\
counts inferred, including a product & 9 & 9 & 11,174 & $3.2\times 10^{4}$ & $1.0\times 10^{7}$ \\
\midrule
\multicolumn{6}{@{}l}{\emph{Other attributes}}\\
needs a schema class the substrate lacks & 24 & 10 & 2,536,941 & $2.6\times 10^{6}$ & $1.1\times 10^{7}$ \\
carries a time axis & 1 & 0 & 1,214,438 & $1.2\times 10^{6}$ & $7.3\times 10^{6}$ \\
figures sourced this pass & 43 & 9 & 166,798,711 & $1.7\times 10^{8}$ & $2.2\times 10^{8}$ \\
figures from recall, to confirm & 7 & 1 & 4,421 & $1.3\times 10^{4}$ & $1.6\times 10^{5}$ \\
instances a reported count & 38 & 9 & 166,759,707 & $1.7\times 10^{8}$ & $2.2\times 10^{8}$ \\
instances a product of the axes & 8 & 0 & 40,425 & $1.6\times 10^{5}$ & $4.0\times 10^{5}$ \\
instances an estimate & 4 & 1 & 3,000 & $3.0\times 10^{3}$ & $3.0\times 10^{3}$ \\
\midrule
Total & 50 & 10 & 166,803,132 & $1.7\times 10^{8}$ & $2.2\times 10^{8}$ \\
\end{longtable}
\endgroup


Table~\ref{tab:summary} gives the pool by host, by band, and by tier, with the record,
genotype, and measurement sums for each. Three properties of those sums are worth stating
before they are read as a ranking.

A handful of rows account for the measurement sums, so the hosts are not comparable on them.
Three \org{S. cerevisiae} rows hold more measurements than every other row in the pool
combined, which makes the yeast total a statement about one platform rather than about the
host, and one artifact accounts for the pseudomonad total, its rows sharing a pipeline and a
condition vocabulary. The new host's total is four orders of magnitude below either, which is
the measured condition of the host rather than a property of the ranking.

Record counts and measurement counts disagree in direction, which is the point of the rank
quantity in Sec.~\ref{sec:measurements}. The isolate panels hold hundreds of records each and
millions of measurements, because one record is a genome plus a transcriptome or a
genome plus a proteome, while the arrayed and pooled screens hold the opposite shape. Summing
either column alone orders the hosts the wrong way round.

And the sums are upper bounds in two specific ways that Sec.~\ref{sec:caveats} records. A
count that is the product of two reported axes assumes a dense matrix, and a real screen drops
strains per condition, so every such figure is an upper bound. A measurement count for a
vector-valued row multiplies by a reported or estimated depth, and a label-free proteome or a
sparse count matrix quantifies a different subset in each sample, so the dense assumption fails
there too.

\subsection{Time\secstatus{todo}}\label{sec:time}

Five rows in this pool carry a time dimension, and they differ in kind as well as in
length. The low-pH expression atlas samples at 0 h and 1 h after the acid shift, so its series
is two points per strain. The 3,053 induction tracks of the fungal language model's corpus
typically carry eight time points per induction, which is the same design as the served
induction series. The longitudinal isolate panel is the longest axis in the pool at 0.2 to 10.2
years of colonization across 52 clone types, where the ordering is part of the phenotype rather
than a label on it. The isoprenol design-build-test-learn ladder is indexed per round across
seven deposits, which is a campaign axis and not a time axis, and it appears here only so that
the two are not conflated. And one allele-specific expression study on the diauxic shift carries 19
time points, the densest series found.

What the schema holds for all of them is \texttt{Environment.duration\_hours} and
\texttt{Environment.duration\_generations}, treatment durations that make two exposure lengths
two distinct environments, and the served induction series was ingested that way as 1,752
records with one environment per time point. That is a flattening: a model sees time as a
categorical condition, and nothing on the record states that eight environments are one series
in order. A first-class time field on the record and a series identity that groups its points
are the two additions a time-resolved model needs, and the five rows above would exercise both.
They are recorded here as an open schema item rather than folded into the
ranking, because the rule of Sec.~\ref{sec:rule} does not score a time axis.
